\documentclass[10pt,a4paper]{amsart}
\usepackage[margin=19mm]{geometry}
\usepackage{amsmath,amssymb,mathtools}
\usepackage[scr=rsfs,cal=euler]{mathalfa}
\usepackage{xfrac}
\usepackage[unicode,hidelinks,bookmarks=false]{hyperref}
\newcommand{\ie}{i.e.\ }
\newcommand{\cf}{cf.\ }
\newcommand{\resp}{respectively\ }
\newcommand{\Subsection}{\subsection}
\providecommand{\nobreakdash}{\nobreak\hbox{-}}
\setlength{\emergencystretch}{2em}
\multlinegap0pt
\usepackage{mdframed}
\usepackage{mdframed}
\usepackage{mdframed}
\usepackage{mdframed}
\usepackage{amssymb}
\usepackage{enumerate}
\usepackage[shortlabels]{enumitem}
\usepackage[normalem]{ulem}
\usepackage{tikz}
\usepackage{mdframed}
\usepackage{xcolor}
\usepackage{float}
\newtheorem{lemma}{Lemma}
\newcommand{\Z}{\mathbb{Z}}
\title{On the Faithful Projective Representations of Finite Groups and their Minimal Dimension}
\author{Sumana Hatui, Poonam Nayak}
\date{Source: arXiv:2605.27280v1 (CC BY 4.0)}
\begin{document}
\maketitle

\noindent\emph{Source and licence.} On the Faithful Projective Representations of Finite Groups and their Minimal Dimension. Version arXiv:2605.27280v1 (CC BY 4.0), distributed by arXiv under the Creative Commons Attribution 4.0 International licence, http://creativecommons.org/licenses/by/4.0/, as recorded in the arXiv licence metadata for this exact version and shown on the abstract page https://arxiv.org/abs/2605.27280v1. Full licence notice: \texttt{licenses/arxiv-2605.27280-CC-BY-4.0.txt}.\par\smallskip

\noindent\textit{Editorial scope.} Counted unit: the article's lemma group (source0.tex lines 816--818). The article's complete original proof of it is delivered with it (source0.tex lines 820--865, one \texttt{proof} environment of 46 lines). No mathematics is written, repaired or re-derived: the statement and the proof are the article's own lines, in the article's own notation, and no retained line is substituted or re-flowed.  No further source-local context is needed: every cross-reference of every retained line resolves inside the two delivered spans.  Nothing was omitted from any retained span, so the package carries no omission record.  No retained line contains a citation, so the wrapper carries no bibliography and no bibliography entry.  No retained line contains a figure environment, an image-inclusion command or drawing code, so no image is delivered and the package declares no asset dependency.  Editorial formatting only, in addition: an AMS \texttt{amsart} preamble that restates the article's own declarations and macros used by the retained lines, namely the environments \texttt{lemma} on separate counters, together with the standard packages those lines need.  The retained environments print in this document as Lemma 1 in order of appearance, whereas the article prints the same items with the numbering of its own full document.  The complete native source of the article as distributed in the arXiv e-print for this exact version is bundled unchanged as \texttt{sources/201483/source0.tex}.
\begin{lemma} \label{1dim direct sum in abelian group}
	Let $G$ be a finite abelian $p$-group of rank $n$ with $|G|>1$. Suppose $\rho= \oplus_{i=1}^r \rho_i$, where $\rho_i \in \mathrm{Irr}(G)$. If $1 \leq r \leq n$, then $\rho$ is not a faithful projective representation of $G$.   
\end{lemma}

\begin{proof}
	The proof is straightforward for $n=1$. So, we assume $n>1$. First, we consider $G=(\mathbb Z/p\mathbb Z)^n$. Let $e_{k}= \underbrace{(0,0, \ldots, 1, \ldots, 0)}_\text{$1$ in the $k{\text {th}}$ position} \in G$ for $1 \leq k \leq n$.
	
	We assume that $r=n$.  
	Suppose $\rho$ is a faithful projective representation of $G$. Then, for each $k$, there exists $j_{l k} \in \{0,1, \ldots, p-1\}$ ($1\leq l \leq n $) such that all $j_{l k}$'s are not equal and $$\rho(e_{k})=
	\left[\begin{array}{ccc}e^{\frac{2 \pi ij_{1k}}{p}} &  & \\
		& \ddots & \\
		& & e^{\frac{2 \pi ij_{nk}}{p}}\end{array}\right]$$
	is a diagonal matrix. Hence, $$\rho\left(m_{1}, m_{2}, \ldots, m_{n}\right)
	=\left[\begin{array}{ccc}e^{\frac{2 \pi i \sum\limits_{k=1}^{n} m_{k} j_{1 k}}{p}} &  & \\
		& \ddots & \\
		& & e^{\frac{2 \pi i \sum \limits_{k=1}^{n} m_{k} j_{n k}}{p}}\end{array}\right].$$                                            
	Since $\rho$ is also a faithful ordinary representation, $\rho\left(m_{1}, m_{2} \ldots, m_{n}\right)= I$ implies that $\left(m_{1}, \ldots, m_{n}\right)= (0,0, \ldots, 0)$, which says
	$$
	\left[\begin{array}{ccc}
		j_{11} & \cdots & j_{1 n} \\
		\vdots & &\vdots \\
		j_{n 1} & \cdots & j_{n n}
	\end{array}\right]\left[\begin{array}{c}
		x_{1} \\
		\vdots \\
		{x}_{n}
	\end{array}\right]= \left[\begin{array}{c}
		0 \\
		\vdots \\
		0
	\end{array}\right]
	$$
	has only a trivial solution in $\Z/p\Z$. Then, the matrix $A=(j_{lk})_{n\times n}$ is invertible,
	and for $b=[1,1,\cdots, 1]^{\top}$, $A x=b$ has a unique solution $c=A^{-1} b,$ where $c \neq 0$.
	Let $c^{\top}= \left[c_{1}, c_{2}, \ldots ., c_{n}\right]$, 
	$c_{i}$ $\in \{0,1, \ldots, p-1\}$.
	Then, $\sum\limits_{k=1}^{n} j_{i k} c_{k}= 1$ for $1\leq i \leq n$. Therefore, 
	\[\rho\left(c_{1}, c_{2}, \ldots, c_{n}\right)=
	\left[\begin{array}{ccc}
		e\frac{2 \pi i \sum\limits_{k=1}^n c_k j_{1k}}{p} & & \\
		& \ddots & \\
		&  & e \frac{2 \pi i \sum\limits_{k=1}^n c_k j_{nk}}{p} 
	\end{array}\right]
	= e^{\frac{2 \pi i}{p}} I_{n \times n},
	\] which contradicts our assumption. Hence, $\rho$ cannot be faithful. 
	
	If $r<n$, let $\psi= \rho \oplus (\oplus_{i=1}^{n-r} 1_G$). If $\rho$ is faithful as projective, then $\psi$ is also faithful as projective, which gives a contradiction to Case 1. Hence, the result holds if $G$ is elementary abelian. 
	Now, let $G$ be any abelian group of rank $n$ and $H= (\Z/p\Z)^n \subseteq G$. If $\rho$ is faithful on $G$, then it is also faithful on $H$, which is not possible. This completes the proof.
	
\end{proof}

\end{document}
